% Grammatikerklaerung nach der hochgeladenen Erklaerungsvorlage

\documentclass[a4paper,11pt]{article}

\usepackage[a4paper,margin=2cm]{geometry}
\usepackage[T1]{fontenc}
\usepackage[utf8]{inputenc}
\usepackage[ngerman,english]{babel}
\usepackage{lmodern}
\usepackage{microtype}
\usepackage[table]{xcolor}
\usepackage{parskip}
\usepackage{enumitem}
\usepackage{array}
\usepackage{longtable}
\usepackage[most]{tcolorbox}
\usepackage{titlesec}
\usepackage[hidelinks]{hyperref}

\definecolor{HeaderBlueDark}{RGB}{70,110,170}
\definecolor{RowBlueLight}{RGB}{235,243,255}
\definecolor{RowBlueDark}{RGB}{220,233,250}
\definecolor{SoftGray}{RGB}{90,90,90}

\setlength{\parindent}{0pt}
\setlist[itemize]{leftmargin=1.4em,itemsep=0.35em,topsep=0.35em}
\renewcommand{\arraystretch}{1.35}
\setlength{\tabcolsep}{5pt}
\linespread{1.06}

\titleformat{\section}[block]
{\Large\bfseries\color{HeaderBlueDark}}
{}{0pt}{}
\titlespacing{\section}{0pt}{1.3em}{0.8em}

\titleformat{\subsection}[block]
{\large\bfseries\color{HeaderBlueDark}}
{}{0pt}{}
\titlespacing{\subsection}{0pt}{1.0em}{0.6em}

\newtcolorbox{exbox}{enhanced,breakable,colback=RowBlueLight,colframe=RowBlueDark,boxrule=0.4pt,arc=1.5mm,left=1.2mm,right=1.2mm,top=1mm,bottom=1mm,title={Beispiele \textnormal{(Examples)}},fonttitle=\bfseries,colbacktitle=RowBlueDark,coltitle=black}

\newtcolorbox{warningbox}{enhanced,colback=RowBlueLight,colframe=HeaderBlueDark,boxrule=0.6pt,arc=1.5mm,left=1.5mm,right=1.5mm,top=1mm,bottom=1mm,title={Achtung \textnormal{(Important)}},fonttitle=\bfseries,colbacktitle=HeaderBlueDark,coltitle=white}

\NewDocumentEnvironment{grammarentry}{mm+b}{%
    \begin{tcolorbox}[enhanced,breakable,colback=white,colframe=HeaderBlueDark,boxrule=0.6pt,arc=1.5mm,left=1.5mm,right=1.5mm,top=1mm,bottom=1mm,colbacktitle=HeaderBlueDark,coltitle=white,fonttitle=\bfseries,title={#1\hfill\normalfont\footnotesize #2}]
    #3
    \end{tcolorbox}
}{}

\newcommand{\GermanText}[1]{\textbf{Deutsch: }#1\par}
\newcommand{\EnglishText}[1]{{\small\color{SoftGray}\textbf{English: }#1\par}}
\newcommand{\ExampleItem}[2]{\item \textbf{#1}\\{\small\color{SoftGray}#2}}

\title{Die Grammatik von \glqq vor\grqq}
\author{}
\date{}

\begin{document}
    \maketitle
    \tableofcontents
    \newpage

    \section{Grundbedeutungen}

        \begin{grammarentry}{vor}{Präposition, Adverb und Verbpräfix}
            \GermanText{Das Wort \textbf{vor} hat mehrere Funktionen. Am wichtigsten ist \textbf{vor} als Präposition. Es kann räumlich \glqq in front of\grqq, zeitlich \glqq before\grqq{} oder \glqq ago\grqq{} und kausal etwa \glqq because of\grqq{} bedeuten.}
            \EnglishText{The word \textbf{vor} has several functions. Most importantly, \textbf{vor} is a preposition. It can mean ``in front of'' spatially, ``before'' or ``ago'' temporally, and approximately ``because of'' when expressing a cause.}

            \GermanText{Außerdem erscheint \textbf{vor} als Adverb und als Bestandteil vieler trennbarer Verben, zum Beispiel \textbf{vorlesen}, \textbf{vorstellen} und \textbf{vorbereiten}.}
            \EnglishText{It also appears as an adverb and as a component of many separable verbs, such as \textbf{vorlesen}, \textbf{vorstellen}, and \textbf{vorbereiten}.}
        \end{grammarentry}

        \begin{exbox}
            \begin{itemize}
                \ExampleItem{Das Auto steht vor dem Haus.}{The car is in front of the house.}
                \ExampleItem{Wir sprechen vor dem Essen.}{We speak before the meal.}
                \ExampleItem{Ich war vor zwei Jahren in Berlin.}{I was in Berlin two years ago.}
                \ExampleItem{Er zittert vor Angst.}{He is trembling with fear.}
            \end{itemize}
        \end{exbox}


    \section{Räumliches \glqq vor\grqq: Dativ oder Akkusativ}

        \begin{grammarentry}{Wechselpräposition}{Wo? = Dativ; Wohin? = Akkusativ}
            \GermanText{Bei einer räumlichen Bedeutung ist \textbf{vor} eine Wechselpräposition. Der Kasus hängt nicht einfach vom Verb ab, sondern davon, ob der Satz eine \textbf{Position} oder eine \textbf{gerichtete Veränderung der Position} beschreibt.}
            \EnglishText{With a spatial meaning, \textbf{vor} is a two-way preposition. The case does not simply depend on the verb, but on whether the sentence describes a \textbf{position} or a \textbf{directed change of position}.}

            \GermanText{Frage \textbf{Wo?}: Dativ. Etwas oder jemand befindet sich bereits an einem Ort vor etwas anderem.}
            \EnglishText{Question \textbf{Where?}: dative. Something or someone is already located in a position in front of something else.}

            \GermanText{Frage \textbf{Wohin?}: Akkusativ. Etwas oder jemand wird an einen Ort vor etwas anderem gebracht oder bewegt.}
            \EnglishText{Question \textbf{Where to?}: accusative. Something or someone is brought or moved to a position in front of something else.}
        \end{grammarentry}

        \rowcolors{2}{RowBlueLight}{RowBlueDark}
        \begin{longtable}{>{\bfseries}p{2.5cm}p{3cm}p{4.8cm}p{4.1cm}}
            \rowcolor{HeaderBlueDark}
            \color{white}\textbf{Bedeutung} &
            \color{white}\textbf{Frage / Kasus} &
            \color{white}\textbf{Beispiel} &
            \color{white}\textbf{English} \\
            Position & Wo? + Dativ & Der Bus wartet vor dem Bahnhof. & The bus is waiting in front of the station. \\
            Richtung / Platzierung & Wohin? + Akkusativ & Der Fahrer stellt den Bus vor den Bahnhof. & The driver positions the bus in front of the station. \\
        \end{longtable}

        \begin{exbox}
            \begin{itemize}
                \ExampleItem{Ich stehe vor der Tür.}{I am standing in front of the door.}
                \ExampleItem{Ich stelle mich vor die Tür.}{I position myself in front of the door.}
                \ExampleItem{Das Fahrrad steht vor dem Gebäude.}{The bicycle is in front of the building.}
                \ExampleItem{Ich stelle das Fahrrad vor das Gebäude.}{I put the bicycle in front of the building.}
            \end{itemize}
        \end{exbox}

        \begin{warningbox}
            \GermanText{Bewegung bedeutet nicht automatisch Akkusativ. Entscheidend ist, ob sich die Position relativ zum Bezugsobjekt ändert. In \glqq Ich gehe vor dem Haus spazieren\grqq{} bleibt die Handlung im Bereich vor dem Haus; deshalb steht der Dativ.}
            \EnglishText{Movement does not automatically require the accusative. What matters is whether the position relative to the reference object changes. In ``Ich gehe vor dem Haus spazieren,'' the action remains in the area in front of the house; therefore, the dative is used.}
        \end{warningbox}


    \section{Artikel und Formen nach \glqq vor\grqq}

        \begin{grammarentry}{Kasusformen}{bestimmter Artikel}
            \GermanText{Bei der räumlichen Verwendung muss man die Dativ- und Akkusativformen der Artikel unterscheiden. Im Plural erhält das Nomen im Dativ normalerweise zusätzlich die Endung \textbf{-n}, sofern es sie nicht schon hat.}
            \EnglishText{With spatial usage, one must distinguish the dative and accusative forms of the articles. In the dative plural, the noun normally also receives the ending \textbf{-n}, unless it already has it.}
        \end{grammarentry}

        \rowcolors{2}{RowBlueLight}{RowBlueDark}
        \begin{longtable}{>{\bfseries}p{2.6cm}p{2.8cm}p{2.8cm}p{2.8cm}p{3.2cm}}
            \rowcolor{HeaderBlueDark}
            \color{white}\textbf{Kasus} &
            \color{white}\textbf{Maskulin} &
            \color{white}\textbf{Feminin} &
            \color{white}\textbf{Neutrum} &
            \color{white}\textbf{Plural} \\
            Dativ & vor dem Mann & vor der Frau & vor dem Kind & vor den Häusern \\
            Akkusativ & vor den Mann & vor die Frau & vor das Kind & vor die Häuser \\
        \end{longtable}

        \begin{exbox}
            \begin{itemize}
                \ExampleItem{Der Hund sitzt vor dem Mann.}{The dog is sitting in front of the man.}
                \ExampleItem{Der Hund läuft vor den Mann.}{The dog runs to a position in front of the man.}
                \ExampleItem{Wir warten vor den Häusern.}{We are waiting in front of the houses.}
            \end{itemize}
        \end{exbox}


    \section{Zeitliches \glqq vor\grqq}

        \subsection{Vor einem Zeitpunkt oder Ereignis}

            \begin{grammarentry}{before}{vor + Dativ}
                \GermanText{Zeitlich bezeichnet \textbf{vor} einen Zeitpunkt, der früher als ein anderer Zeitpunkt oder ein Ereignis liegt. In dieser Verwendung steht nach \textbf{vor} der Dativ.}
                \EnglishText{Temporally, \textbf{vor} refers to a point in time that is earlier than another point or event. In this use, \textbf{vor} takes the dative.}
            \end{grammarentry}

            \begin{exbox}
                \begin{itemize}
                    \ExampleItem{Vor dem Unterricht trinke ich Kaffee.}{I drink coffee before class.}
                    \ExampleItem{Bitte ruf mich vor der Besprechung an.}{Please call me before the meeting.}
                    \ExampleItem{Kurz vor Mitternacht ging er nach Hause.}{He went home shortly before midnight.}
                    \ExampleItem{Vor Beginn des Kurses müssen Sie sich anmelden.}{You must register before the course begins.}
                \end{itemize}
            \end{exbox}

        \subsection{Eine Zeitspanne zurück: \glqq ago\grqq}

            \begin{grammarentry}{ago}{vor + Zeitangabe}
                \GermanText{Mit einer Zeitspanne bezeichnet \textbf{vor}, wie lange ein vergangenes Ereignis zurückliegt. Die Zeitspanne steht im Dativ, aber ohne Artikel ist der Kasus an manchen Formen nicht sichtbar.}
                \EnglishText{With a period of time, \textbf{vor} indicates how long ago a past event occurred. The time expression is in the dative, although without an article the case is not visible in some forms.}

                \GermanText{Typische Struktur: \textbf{vor + Zahl + Zeiteinheit}.}
                \EnglishText{Typical structure: \textbf{vor + number + unit of time}.}
            \end{grammarentry}

            \begin{exbox}
                \begin{itemize}
                    \ExampleItem{Ich bin vor drei Tagen angekommen.}{I arrived three days ago.}
                    \ExampleItem{Sie hat vor einem Monat eine neue Arbeit begonnen.}{She started a new job one month ago.}
                    \ExampleItem{Vor vielen Jahren wohnte ich in dieser Stadt.}{Many years ago, I lived in this city.}
                \end{itemize}
            \end{exbox}

        \subsection{Uhrzeit}

            \begin{grammarentry}{fünf vor acht}{Angabe einer Uhrzeit}
                \GermanText{Bei Uhrzeiten bedeutet \textbf{vor}, dass eine bestimmte Zahl von Minuten bis zur genannten Stunde fehlt.}
                \EnglishText{With clock times, \textbf{vor} means that a certain number of minutes remain until the stated hour.}
            \end{grammarentry}

            \begin{exbox}
                \begin{itemize}
                    \ExampleItem{Es ist fünf vor acht.}{It is five minutes to eight.}
                    \ExampleItem{Der Zug fährt um Viertel vor neun ab.}{The train departs at a quarter to nine.}
                \end{itemize}
            \end{exbox}


    \section{\glqq vor\grqq und \glqq seit\grqq}

        \begin{grammarentry}{Vergangenheit}{ein Zeitpunkt oder eine Dauer}
            \GermanText{\textbf{vor} nennt den Abstand zu einem vergangenen Zeitpunkt. Die Handlung kann beendet sein.}
            \EnglishText{\textbf{vor} states the distance to a point in the past. The action may be finished.}

            \GermanText{\textbf{seit} bezeichnet den Beginn einer Situation, die bis zur Gegenwart andauert. Nach \textbf{seit} steht ebenfalls der Dativ.}
            \EnglishText{\textbf{seit} marks the beginning of a situation that continues up to the present. \textbf{seit} also takes the dative.}
        \end{grammarentry}

        \begin{exbox}
            \begin{itemize}
                \ExampleItem{Ich bin vor zwei Jahren nach Wien gezogen.}{I moved to Vienna two years ago.}
                \ExampleItem{Ich wohne seit zwei Jahren in Wien.}{I have been living in Vienna for two years.}
                \ExampleItem{Wir haben uns vor einer Woche getroffen.}{We met one week ago.}
                \ExampleItem{Wir kennen uns seit einer Woche.}{We have known each other for one week.}
            \end{itemize}
        \end{exbox}


    \section{Kausales \glqq vor\grqq}

        \begin{grammarentry}{Ursache einer Reaktion}{vor + Dativ}
            \GermanText{\textbf{vor} kann den Grund einer körperlichen, emotionalen oder unwillkürlichen Reaktion ausdrücken. Auch hier steht der Dativ.}
            \EnglishText{\textbf{vor} can express the cause of a physical, emotional, or involuntary reaction. The dative is also used here.}

            \GermanText{Häufige Verbindungen sind \textbf{vor Angst}, \textbf{vor Freude}, \textbf{vor Kälte}, \textbf{vor Hunger}, \textbf{vor Müdigkeit} und \textbf{vor Schmerz}.}
            \EnglishText{Common combinations include \textbf{vor Angst} (from fear), \textbf{vor Freude} (with joy), \textbf{vor Kälte} (from cold), \textbf{vor Hunger} (from hunger), \textbf{vor Müdigkeit} (from tiredness), and \textbf{vor Schmerz} (from pain).}
        \end{grammarentry}

        \begin{exbox}
            \begin{itemize}
                \ExampleItem{Das Kind weinte vor Angst.}{The child cried with fear.}
                \ExampleItem{Sie sprang vor Freude in die Luft.}{She jumped into the air with joy.}
                \ExampleItem{Meine Hände zittern vor Kälte.}{My hands are trembling from the cold.}
                \ExampleItem{Er konnte vor Müdigkeit kaum sprechen.}{He could hardly speak because of tiredness.}
            \end{itemize}
        \end{exbox}

        \begin{warningbox}
            \GermanText{Bei einer absichtlichen Begründung verwendet man normalerweise nicht \textbf{vor}, sondern zum Beispiel \textbf{wegen}, \textbf{aus} oder einen Nebensatz mit \textbf{weil}. \glqq Ich lerne vor der Prüfung\grqq{} bedeutet zeitlich \glqq before the exam\grqq, nicht \glqq because of the exam\grqq.}
            \EnglishText{For an intentional reason, German normally does not use \textbf{vor}, but for example \textbf{wegen}, \textbf{aus}, or a subordinate clause with \textbf{weil}. ``Ich lerne vor der Prüfung'' means temporally ``before the exam,'' not ``because of the exam.''}
        \end{warningbox}


    \section{Verben, Adjektive und Nomen mit \glqq vor\grqq}

        \begin{grammarentry}{Feste Rektion}{vor + Dativ}
            \GermanText{Manche Verben, Adjektive und Nomen verlangen eine Ergänzung mit \textbf{vor + Dativ}. In solchen Fällen muss die Verbindung als Einheit gelernt werden.}
            \EnglishText{Some verbs, adjectives, and nouns require a complement with \textbf{vor + dative}. In such cases, the combination should be learned as a unit.}
        \end{grammarentry}

        \rowcolors{2}{RowBlueLight}{RowBlueDark}
        \begin{longtable}{>{\bfseries}p{4cm}p{5.5cm}p{5.3cm}}
            \rowcolor{HeaderBlueDark}
            \color{white}\textbf{Verbindung} &
            \color{white}\textbf{Beispiel} &
            \color{white}\textbf{English} \\
            Angst haben vor & Ich habe Angst vor der Prüfung. & I am afraid of the exam. \\
            sich fürchten vor & Das Kind fürchtet sich vor dem Hund. & The child is afraid of the dog. \\
            warnen vor & Sie warnt uns vor der Gefahr. & She warns us about the danger. \\
            schützen vor & Die Jacke schützt mich vor der Kälte. & The jacket protects me from the cold. \\
            fliehen vor & Viele Menschen fliehen vor dem Krieg. & Many people flee from the war. \\
            sich verstecken vor & Er versteckt sich vor seinem Chef. & He is hiding from his boss. \\
            sicher vor & Die Daten sind vor unbefugtem Zugriff sicher. & The data is safe from unauthorized access. \\
            Respekt vor & Ich habe großen Respekt vor ihrer Arbeit. & I have great respect for her work. \\
            Vorrang vor & Sicherheit hat Vorrang vor Geschwindigkeit. & Safety takes priority over speed. \\
        \end{longtable}


    \section{Personen, Sachen und Pronominaladverbien}

        \begin{grammarentry}{vor ihm oder davor}{Person versus Sache}
            \GermanText{Bei Personen verwendet man \textbf{vor + Personalpronomen}: \textbf{vor mir}, \textbf{vor dir}, \textbf{vor ihm}, \textbf{vor ihr}, \textbf{vor uns}, \textbf{vor euch}, \textbf{vor ihnen}.}
            \EnglishText{With people, German uses \textbf{vor + personal pronoun}: \textbf{vor mir}, \textbf{vor dir}, \textbf{vor ihm}, \textbf{vor ihr}, \textbf{vor uns}, \textbf{vor euch}, \textbf{vor ihnen}.}

            \GermanText{Bei Sachen, Situationen oder ganzen Aussagen verwendet man häufig das Pronominaladverb \textbf{davor}. Die Frageform ist \textbf{wovor?}. Da \textbf{vor} mit einem Konsonanten beginnt, heißt es \textbf{davor} und \textbf{wovor}, nicht \glqq darvor\grqq{} oder \glqq worvor\grqq.}
            \EnglishText{With things, situations, or entire statements, German often uses the pronominal adverb \textbf{davor}. The question form is \textbf{wovor?}. Because \textbf{vor} begins with a consonant, the forms are \textbf{davor} and \textbf{wovor}, not ``darvor'' or ``worvor.''}
        \end{grammarentry}

        \begin{exbox}
            \begin{itemize}
                \ExampleItem{Ich habe Angst vor meinem Nachbarn. Ich habe Angst vor ihm.}{I am afraid of my neighbor. I am afraid of him.}
                \ExampleItem{Ich habe Angst vor der Prüfung. Ich habe Angst davor.}{I am afraid of the exam. I am afraid of it.}
                \ExampleItem{Wovor hast du Angst?}{What are you afraid of?}
                \ExampleItem{Ich warne dich davor, allein dorthin zu gehen.}{I warn you against going there alone.}
                \ExampleItem{Sie hat Angst davor, dass sie die Prüfung nicht besteht.}{She is afraid that she will not pass the exam.}
            \end{itemize}
        \end{exbox}

        \begin{grammarentry}{davor, dass / davor, zu}{Nebensatz und Infinitivgruppe}
            \GermanText{Eine Ergänzung kann durch \textbf{davor, dass + Nebensatz} ausgedrückt werden. Wenn das Subjekt beider Handlungen gleich ist, ist oft auch \textbf{davor, zu + Infinitiv} möglich.}
            \EnglishText{A complement can be expressed with \textbf{davor, dass + subordinate clause}. When the subject of both actions is the same, \textbf{davor, zu + infinitive} is often also possible.}
        \end{grammarentry}

        \begin{exbox}
            \begin{itemize}
                \ExampleItem{Er fürchtet sich davor, dass er seinen Arbeitsplatz verliert.}{He is afraid that he will lose his job.}
                \ExampleItem{Er fürchtet sich davor, seinen Arbeitsplatz zu verlieren.}{He is afraid of losing his job.}
            \end{itemize}
        \end{exbox}


    \section{\glqq vor\grqq, \glqq bevor\grqq und \glqq vorher\grqq}

        \begin{grammarentry}{Drei verschiedene Wortarten}{Präposition, Konjunktion, Adverb}
            \GermanText{\textbf{vor} ist hier eine Präposition und steht vor einem Nomen oder Pronomen: \textbf{vor dem Essen}.}
            \EnglishText{Here, \textbf{vor} is a preposition and stands before a noun or pronoun: \textbf{vor dem Essen} (before the meal).}

            \GermanText{\textbf{bevor} ist eine unterordnende Konjunktion. Sie leitet einen Nebensatz ein; das konjugierte Verb steht am Ende: \textbf{bevor ich esse}.}
            \EnglishText{\textbf{bevor} is a subordinating conjunction. It introduces a subordinate clause; the conjugated verb stands at the end: \textbf{bevor ich esse} (before I eat).}

            \GermanText{\textbf{vorher} ist ein Adverb und kann ohne nachfolgendes Nomen stehen: \textbf{Ich esse vorher}.}
            \EnglishText{\textbf{vorher} is an adverb and can stand without a following noun: \textbf{Ich esse vorher} (I eat beforehand).}
        \end{grammarentry}

        \rowcolors{2}{RowBlueLight}{RowBlueDark}
        \begin{longtable}{>{\bfseries}p{2.6cm}p{3.6cm}p{4.7cm}p{4.1cm}}
            \rowcolor{HeaderBlueDark}
            \color{white}\textbf{Wort} &
            \color{white}\textbf{Wortart} &
            \color{white}\textbf{Beispiel} &
            \color{white}\textbf{English} \\
            vor & Präposition & Vor dem Schlafengehen lese ich. & I read before going to sleep. \\
            bevor & Konjunktion & Bevor ich schlafen gehe, lese ich. & Before I go to sleep, I read. \\
            vorher & Adverb & Ich lese vorher. & I read beforehand. \\
        \end{longtable}


    \section{Wortstellung mit einer \glqq vor\grqq-Gruppe}

        \begin{grammarentry}{Satzglied}{Position im Hauptsatz}
            \GermanText{Eine Präpositionalgruppe mit \textbf{vor} ist ein Satzglied und kann oft verschoben werden. Steht sie im Vorfeld, folgt im deutschen Hauptsatz direkt das konjugierte Verb. Das Subjekt steht dann normalerweise nach dem Verb.}
            \EnglishText{A prepositional phrase with \textbf{vor} is a sentence element and can often be moved. When it occupies the first position, the conjugated verb follows directly in a German main clause. The subject then normally follows the verb.}
        \end{grammarentry}

        \begin{exbox}
            \begin{itemize}
                \ExampleItem{Ein Auto steht vor dem Haus.}{A car is standing in front of the house.}
                \ExampleItem{Vor dem Haus steht ein Auto.}{In front of the house, there is a car.}
                \ExampleItem{Ich trinke vor der Arbeit einen Kaffee.}{I drink a coffee before work.}
                \ExampleItem{Vor der Arbeit trinke ich einen Kaffee.}{Before work, I drink a coffee.}
            \end{itemize}
        \end{exbox}

        \begin{warningbox}
            \GermanText{Falsch ist: \glqq Vor der Arbeit ich trinke einen Kaffee.\grqq{} Nach dem ersten Satzglied muss im Hauptsatz das konjugierte Verb an zweiter Stelle stehen.}
            \EnglishText{Incorrect: ``Vor der Arbeit ich trinke einen Kaffee.'' After the first sentence element, the conjugated verb must be in second position in a main clause.}
        \end{warningbox}


    \section{Zusammenziehungen: \glqq vorm\grqq und \glqq vors\grqq}

        \begin{grammarentry}{Kurzformen}{vor dem / vor das}
            \GermanText{\textbf{vorm} ist die Zusammenziehung von \textbf{vor dem}. Sie ist besonders in der gesprochenen Sprache häufig und kommt auch in normaler geschriebener Sprache vor.}
            \EnglishText{\textbf{vorm} is the contraction of \textbf{vor dem}. It is especially common in spoken language and also occurs in ordinary written language.}

            \GermanText{\textbf{vors} ist die Zusammenziehung von \textbf{vor das}. Sie wird ebenfalls verwendet, wirkt aber stärker umgangssprachlich. In formellen Texten kann man die vollständige Form schreiben.}
            \EnglishText{\textbf{vors} is the contraction of \textbf{vor das}. It is also used, but sounds more colloquial. In formal texts, the full form can be written.}
        \end{grammarentry}

        \begin{exbox}
            \begin{itemize}
                \ExampleItem{Wir treffen uns vorm Kino.}{We are meeting in front of the cinema.}
                \ExampleItem{Stell den Stuhl vors Fenster.}{Put the chair in front of the window.}
            \end{itemize}
        \end{exbox}


    \section{Feste Ausdrücke mit \glqq vor\grqq}

        \begin{grammarentry}{Häufige Wendungen}{Bedeutung aus dem Zusammenhang}
            \GermanText{Einige häufige Wendungen müssen als feste Einheiten gelernt werden. Ihre Bedeutung lässt sich nicht immer nur aus der räumlichen Grundbedeutung erklären.}
            \EnglishText{Some common expressions should be learned as fixed units. Their meaning cannot always be explained solely from the basic spatial meaning.}
        \end{grammarentry}

        \rowcolors{2}{RowBlueLight}{RowBlueDark}
        \begin{longtable}{>{\bfseries}p{3.8cm}p{5.3cm}p{5.7cm}}
            \rowcolor{HeaderBlueDark}
            \color{white}\textbf{Ausdruck} &
            \color{white}\textbf{Bedeutung / Beispiel} &
            \color{white}\textbf{English} \\
            vor allem & hauptsächlich: Vor allem brauche ich Zeit. & above all; mainly: Above all, I need time. \\
            nach wie vor & weiterhin, noch immer: Das Problem besteht nach wie vor. & still; as before: The problem still exists. \\
            vor Ort & am betreffenden Ort: Ein Techniker ist vor Ort. & on site: A technician is on site. \\
            vor sich hin & ohne bestimmtes Ziel oder ohne andere zu beachten: Er redet vor sich hin. & to oneself / aimlessly: He is talking to himself. \\
            etwas vor Augen haben & sich etwas deutlich vorstellen: Ich habe das Ziel vor Augen. & to have something clearly in mind. \\
            jemandem etwas vor Augen führen & etwas deutlich zeigen oder erklären & to make something clear to someone. \\
            vor Kurzem & vor kurzer Zeit: Ich habe ihn vor Kurzem gesehen. & recently: I saw him recently. \\
        \end{longtable}


    \section{\glqq vor\grqq als Verbpräfix}

        \begin{grammarentry}{vor-}{meist betont und trennbar}
            \GermanText{Bei vielen Verben ist \textbf{vor-} ein betontes, trennbares Präfix. Im Hauptsatz steht der abgetrennte Präfixteil am Satzende. Im Infinitiv und im Nebensatz bleibt das Verb zusammengeschrieben. Im Partizip II steht häufig \textbf{ge} zwischen Präfix und Verbstamm.}
            \EnglishText{In many verbs, \textbf{vor-} is a stressed, separable prefix. In a main clause, the separated prefix appears at the end of the sentence. In the infinitive and in a subordinate clause, the verb remains written as one word. In the past participle, \textbf{ge} often stands between the prefix and the verb stem.}
        \end{grammarentry}

        \rowcolors{2}{RowBlueLight}{RowBlueDark}
        \begin{longtable}{>{\bfseries}p{3.1cm}p{4.5cm}p{4.5cm}p{3.1cm}}
            \rowcolor{HeaderBlueDark}
            \color{white}\textbf{Verb} &
            \color{white}\textbf{Hauptsatz} &
            \color{white}\textbf{Nebensatz} &
            \color{white}\textbf{Partizip II} \\
            vorlesen & Ich lese den Text vor. & ..., weil ich den Text vorlese. & vorgelesen \\
            vorbereiten & Ich bereite alles vor. & ..., weil ich alles vorbereite. & vorbereitet \\
            vorstellen & Ich stelle mich vor. & ..., weil ich mich vorstelle. & vorgestellt \\
            vorhaben & Ich habe heute viel vor. & ..., weil ich heute viel vorhabe. & vorgehabt \\
            vorkommen & Das kommt oft vor. & ..., weil das oft vorkommt. & vorgekommen \\
        \end{longtable}

        \begin{warningbox}
            \GermanText{Nicht jedes längere Verb, das mit den Buchstaben \textbf{vor} beginnt, folgt automatisch genau demselben Muster. Bei neuen Verben sollte man deshalb im Wörterbuch prüfen, ob das Verb trennbar ist und wie sein Partizip II lautet.}
            \EnglishText{Not every longer verb beginning with the letters \textbf{vor} automatically follows exactly the same pattern. Therefore, with new verbs, one should check in a dictionary whether the verb is separable and how its past participle is formed.}
        \end{warningbox}


    \section{Typische Fehler}

        \begin{grammarentry}{Fehler 1}{falscher Kasus bei einer Position}
            \GermanText{Falsch: \glqq Das Auto steht vor das Haus.\grqq{} Richtig: \glqq Das Auto steht vor dem Haus.\grqq{} Das Auto befindet sich an einer Position; deshalb steht der Dativ.}
            \EnglishText{Incorrect: ``Das Auto steht vor das Haus.'' Correct: ``Das Auto steht vor dem Haus.'' The car is in a position; therefore, the dative is used.}
        \end{grammarentry}

        \begin{grammarentry}{Fehler 2}{falscher Kasus bei einer Platzierung}
            \GermanText{Falsch: \glqq Ich stelle das Auto vor dem Haus.\grqq{} Richtig, wenn eine neue Position gemeint ist: \glqq Ich stelle das Auto vor das Haus.\grqq{} Die Handlung beantwortet die Frage \glqq Wohin?\grqq.}
            \EnglishText{Incorrect: ``Ich stelle das Auto vor dem Haus.'' Correct, when a new position is meant: ``Ich stelle das Auto vor das Haus.'' The action answers the question ``Where to?''}
        \end{grammarentry}

        \begin{grammarentry}{Fehler 3}{\glqq vor\grqq statt \glqq bevor\grqq}
            \GermanText{Falsch: \glqq Vor ich gehe, rufe ich dich an.\grqq{} Richtig: \glqq Bevor ich gehe, rufe ich dich an.\grqq{} Vor einem vollständigen Nebensatz braucht man die Konjunktion \textbf{bevor}.}
            \EnglishText{Incorrect: ``Vor ich gehe, rufe ich dich an.'' Correct: ``Bevor ich gehe, rufe ich dich an.'' Before a complete subordinate clause, the conjunction \textbf{bevor} is required.}
        \end{grammarentry}

        \begin{grammarentry}{Fehler 4}{\glqq vor\grqq und \glqq seit\grqq verwechselt}
            \GermanText{Falsch für eine noch andauernde Situation: \glqq Ich wohne vor zwei Jahren hier.\grqq{} Richtig: \glqq Ich wohne seit zwei Jahren hier.\grqq{} Oder für ein abgeschlossenes Ereignis: \glqq Ich bin vor zwei Jahren hierhergezogen.\grqq}
            \EnglishText{Incorrect for a situation that is still continuing: ``Ich wohne vor zwei Jahren hier.'' Correct: ``Ich wohne seit zwei Jahren hier.'' Or for a completed event: ``Ich bin vor zwei Jahren hierhergezogen.''}
        \end{grammarentry}

        \begin{grammarentry}{Fehler 5}{Personen durch \glqq davor\grqq ersetzt}
            \GermanText{Bei Personen sagt man normalerweise \textbf{vor ihm}, \textbf{vor ihr} oder \textbf{vor ihnen}, nicht \textbf{davor}.}
            \EnglishText{With people, one normally says \textbf{vor ihm}, \textbf{vor ihr}, or \textbf{vor ihnen}, not \textbf{davor}.}
        \end{grammarentry}


    \newpage
    \section{Entscheidungsschema}

        \begin{grammarentry}{So wählst du die richtige Form}{Schritt für Schritt}
            \GermanText{1. Hat \textbf{vor} eine räumliche Bedeutung? Frage dann, ob eine Position oder eine gerichtete Platzierung gemeint ist.}
            \EnglishText{1. Does \textbf{vor} have a spatial meaning? Then ask whether a position or a directed placement is meant.}

            \GermanText{2. Position, Frage \glqq Wo?\grqq: Dativ. Gerichtete Platzierung, Frage \glqq Wohin?\grqq: Akkusativ.}
            \EnglishText{2. Position, question ``Where?'': dative. Directed placement, question ``Where to?'': accusative.}

            \GermanText{3. Hat \textbf{vor} eine zeitliche, kausale oder fest regierte Bedeutung? Dann steht normalerweise der Dativ.}
            \EnglishText{3. Does \textbf{vor} have a temporal, causal, or lexically governed meaning? Then the dative is normally used.}

            \GermanText{4. Folgt ein vollständiger Nebensatz? Dann brauchst du häufig \textbf{bevor} oder eine Verbindung mit \textbf{davor, dass}.}
            \EnglishText{4. Is a complete subordinate clause following? Then you often need \textbf{bevor} or a construction with \textbf{davor, dass}.}

            \GermanText{5. Bezieht sich die Ergänzung auf eine Sache? Dann kann oft \textbf{davor} stehen. Bezieht sie sich auf eine Person? Dann verwende \textbf{vor + Personalpronomen}.}
            \EnglishText{5. Does the complement refer to a thing? Then \textbf{davor} can often be used. Does it refer to a person? Then use \textbf{vor + personal pronoun}.}
        \end{grammarentry}


    \section{Zusammenfassung}

        \begin{grammarentry}{Merksätze}{kurz und wichtig}
            \GermanText{Räumlich: \textbf{Wo? + Dativ}; \textbf{Wohin? + Akkusativ}.}
            \EnglishText{Spatially: \textbf{Where? + dative}; \textbf{Where to? + accusative}.}

            \GermanText{Zeitlich: \textbf{vor + Dativ} = before oder ago.}
            \EnglishText{Temporally: \textbf{vor + dative} = before or ago.}

            \GermanText{Kausal und in festen Verbindungen: normalerweise \textbf{vor + Dativ}.}
            \EnglishText{Causally and in fixed combinations: normally \textbf{vor + dative}.}

            \GermanText{Vor einem Nomen: \textbf{vor}; vor einem Nebensatz: häufig \textbf{bevor}; allein als Zeitadverb: \textbf{vorher}.}
            \EnglishText{Before a noun: \textbf{vor}; before a subordinate clause: often \textbf{bevor}; alone as a temporal adverb: \textbf{vorher}.}

            \GermanText{Sache: häufig \textbf{davor}; Person: \textbf{vor ihm / ihr / ihnen}.}
            \EnglishText{Thing: often \textbf{davor}; person: \textbf{vor ihm / ihr / ihnen}.}
        \end{grammarentry}

\end{document}
